% 4.1 第六章 电荷与电场.tex
\chapter{第六章 电荷与电场}

\section{库仑定律与电场强度}
    \subsection{电荷及其性质}
    \subsection{库仑定律}
    \[
    F=k\frac{q_1q_2}{r^3}\boldsymbol{r}
    \]
    比例系数$k$写成
    \[
    k=\frac{1}{4\pi\varepsilon_0}
    \]
    式中，$\varepsilon_0$称为真空介电常量，也叫真空电容率。
    %\[
    %\varepsilon_0\approx 8.85\times10^{-12}\,\mathrm{F/m}
    %\]
    \subsection{电场与电场强度}
    \subsection{场强叠加原理}

\section{电场强度通量与高斯定理}
    \subsection{电场线}
    \subsection{电场强度通量}
    通过某个闭合曲面$S$的$\boldsymbol{E}$通量
    \subsection{真空中高斯定理}
    在真空中静电场内通过任意闭合曲面的$\boldsymbol{E}$通量，等于该闭合曲面所包围的电荷的代数和除以，其数学表达式为
    %\[
    %\]
    若场源电荷连续分布，则式为
    %\[
    %\]
    \subsection{高斯定理的应用}

\section{静电场的环路定理与电势}
    \subsection{静电场的环路定理}
    在静电场中，场强沿闭合路径的线积分等于零，称静电场环路定理。
    \subsection{电势}
    \subsection{电势差}
    \subsection{电势的计算}
    \subsection{等势面 电势梯度}
        \subsubsection{等势面}
        \subsubsection{电势梯度}

\section{静电场中导体}
    \subsection{导体的静电平衡性质}
        \subsubsection{导体的静电平衡状态}
        \subsubsection{静电平衡下导体带电特性}
    \subsection{空腔导体和静电屏蔽}
        \subsubsection{空腔导体}
        \subsubsection{静电屏蔽}
    \subsection{电容和电容器}
        \subsubsection{孤立导体的电容}
        \subsubsection{电容器及其电容}
        \subsubsection{串联和并联}
            \subsubsubsection{串联}
            \subsubsubsection{并联}

\section{静电场中的电介质 电介质中的高斯定理}
    \subsection{电介质的微观结构和极化过程}
    \subsection{电极化强度矢量}
    \subsection{电介质中的场强}
    \subsection{电介质中的高斯定理 电位移\texorpdfstring{$\boldsymbol{D}$}{D}}
    \subsection{有电介质时静电场的计算}

\section{电场的能量}
    \subsection{点电荷系统的电能}
    \subsection{电容器的能量}
    \subsection{电场能量 电场能量密度}
    \subsection{电场能量的计算}
    电场能量为
    %\[
    %d
    %\]
    即可求得整个电场的总能量为
    %\[
    %W_e
    %\]
    显然

\endinput